% Terjemahan Bahasa Indonesia dari prelude.tex baris 30-103.
% Sumber: Methods of Algebra, Volume 2, commit
% 9a5803ff77dd3257484cb177f851a73770a59dd3 (CC BY 4.0).
% stable-unit-id: o014.aljabr2.prelude.homology-history

\section*{Sejarah Singkat Homologi}
\addcontentsline{toc}{section}{Sejarah Singkat Homologi}
\hypertarget{o014.aljabr2.prelude.homology-history}{ }

% segment-id: o014.li2.prelude.homology-history.p001
Matematika tidak harus berpijak pada sejarah. Namun, menelaah masa lalu
untuk memahami masa depan tentu bermanfaat dan tidak merugikan. Untuk memudahkan pembahasan,
kita akan membagi perkembangan teori homologi secara agak artifisial ke
dalam beberapa tahap.

\subsection*{Masa Perintisan}

% segment-id: o014.li2.prelude.homology-history.p002
Pendorong pertama perkembangan aljabar homologis adalah topologi. Pada
paruh akhir abad ke-19, B.\ Riemann dan E.\ Betti mengkaji genus permukaan
beserta generalisasinya ke dimensi lebih tinggi, yang disebut bilangan
Betti. Menjelang pergantian abad, gagasan ini dikembangkan oleh
H.\ Poincaré menjadi suatu sistem yang lebih ketat secara matematis. Secara kasar, gagasan
awal Poincaré ialah menguraikan ruang menjadi poligon, polihedron, atau
padanan berdimensi lebih tinggi yang direkatkan satu sama lain, lalu
menghitung bilangan Betti atau versi torsinya dari matriks yang merekam
cara perekatan tersebut, dan menunjukkan bahwa semuanya merupakan
invarian topologis.

% segment-id: o014.li2.prelude.homology-history.p003
Dalam tulisan singkat tahun 1925, E.\ Noether menjelaskan
bagaimana homologi dapat didefinisikan sebagai grup abelian, sedangkan
bilangan Betti dan versi torsinya hanyalah nilai numerik yang diturunkan
darinya. Gagasan ini sudah mendekati definisi modern dalam topologi
aljabar: untuk ruang $E$ yang dilengkapi dengan suatu penguraian dan grup
abelian $A$, didefinisikan
kompleks rantai $C(E, A)$ yang merekam cara bagian-bagian penguraian
tersebut direkatkan.
Grup homologi $E$ berkoefisien $A$ kemudian didefinisikan sebagai
$\Hm_n(E; A) := \Hm_n(C(E, A))$, sementara grup kohomologi
$\Hm^n(E; A)$ didefinisikan melalui cara yang dual.

% segment-id: o014.li2.prelude.homology-history.p004
Sejak saat itu hingga sekitar tahun 1950, topologi aljabar memasuki masa
perkembangan yang amat subur. Dari sudut pandang aljabar, perkembangan
yang patut dicatat meliputi:
\begin{itemize}
	% segment-id: o014.li2.prelude.homology-history.p005
	\item Teorema koefisien universal untuk homologi (atau kohomologi),
	yang menyatakan bahwa kasus berkoefisien bilangan bulat sudah cukup
	untuk menentukan homologi (atau kohomologi) dengan koefisien dalam grup
	abelian umum $A$. Uraian konkretnya melibatkan funktor $\Tor$ (atau
	$\Ext$), yang akan dibahas secara mendalam kelak.
	% segment-id: o014.li2.prelude.homology-history.p006
	\item Hasil kali \emph{cup} pada kohomologi, yaitu suatu struktur
	perkalian yang dimiliki grup kohomologi.
	% segment-id: o014.li2.prelude.homology-history.p007
	\item Ruang yang disebut asferis, yakni ruang terhubung yang memenuhi
	$n > 1 \implies \pi_n(E) = 0$; di sini $E$ menyatakan ruang tersebut,
	dan $\pi_n(E)$ menyatakan grup homotopi ke-$n$ relatif terhadap titik
	dasar yang dipilih.
	W.\ Hurewicz membuktikan bahwa grup homologi dan kohomologi ruang
	asferis sepenuhnya ditentukan oleh $\pi_1(E)$ dan dapat diberikan secara
	aljabar. Inilah cikal bakal teori kohomologi grup.
\end{itemize}

% segment-id: o014.li2.prelude.homology-history.p008
Pendorong lain bagi aljabar homologis datang dari dalam aljabar sendiri.
Contoh klasiknya ialah teorema sizigi Hilbert dalam teori invarian (1890):
dalam bahasa modern, teorema sizigi menyatakan bahwa setiap modul atas
gelanggang polinomial $\Bbbk[X_1, \ldots, X_n]$ memiliki resolusi bebas
dengan panjang $\leq n$.

% segment-id: o014.li2.prelude.homology-history.p009
Funktor $\Ext$ juga berkaitan dengan persoalan aljabar. Untuk grup abelian
$A$ dan $B$, barisan eksak pendek grup abelian
$0 \to A \to E \to B \to 0$ disebut pula ekstensi grup (perluasan grup);
hingga isomorfisme, ekstensi-ekstensi tersebut berkorespondensi satu-satu
dengan unsur-unsur grup $\Ext^1(B, A)$. Selain itu, jika $B$ diganti oleh
sembarang grup $G$, klasifikasi ekstensi grup secara alami menghasilkan
definisi aljabar bagi kohomologi grup $\Hm^2(G, A)$.

% segment-id: o014.li2.prelude.homology-history.p010
Operasi derivasi pada gelanggang merupakan persoalan yang sejak awal
menarik perhatian para ahli aljabar. Pada tahun 1942, G.\ Hochschild
mendefinisikan invarian untuk gelanggang $R$ yang kini disebut homologi
dan kohomologi Hochschild; suku kohomologi berderajat $n=1$
mengklasifikasikan derivasi modulo derivasi dalam.\footnote{Catatan
penerjemah: teks sumber menyatakan bahwa suku berderajat satu
``mengklasifikasikan semua derivasi''. Pernyataan ini dikoreksi di sini:
$\Hm^1_{\mathrm{Hoch}}(R,R)$ mengklasifikasikan derivasi modulo derivasi
dalam; lihat catatan koreksi O014-C001.}
Konstruksi-konstruksi ini dapat diperluas ke bimodul $(R, R)$ yang lebih
umum, yaitu $M$, dan terus memainkan peranan penting dalam matematika
mutakhir.

% segment-id: o014.li2.prelude.homology-history.p011
Patut disebut pula karya J.\ Leray mengenai berkas dan barisan spektral
selama Perang Dunia II. Motivasinya bermula dari persoalan titik tetap
yang berkaitan dengan persamaan diferensial parsial. Berkas adalah
struktur dalam geometri yang memuat informasi lokal--global, sedangkan
barisan spektral menyediakan perangkat ampuh untuk menghitung kohomologi
berkas. Di bawah pengaruh teori Leray dan karya Kiyoshi Oka dalam teori
fungsi beberapa peubah kompleks, H.\ Cartan dan sejumlah matematikawan lain
mulai menulis ulang dasar-dasar topologi aljabar setelah perang; aljabar
homologis pun memperoleh wajah yang sama sekali baru.

\subsection*{Revolusi Cartan--Eilenberg}

% segment-id: o014.li2.prelude.homology-history.p012
Karya H.\ Cartan dan S.\ Eilenberg \cite{CE56} sangat luas dan mendalam,
serta untuk waktu yang lama dipandang sebagai rujukan kanonik. Buku itu
menandai suatu masa baru dalam perkembangan aljabar homologis. Sumbangan
utamanya mencakup:
\begin{itemize}
	% segment-id: o014.li2.prelude.homology-history.p013
	\item definisi modul injektif dan modul projektif, beserta resolusi
	injektif dan resolusi projektif suatu modul;
	% segment-id: o014.li2.prelude.homology-history.p014
	\item definisi funktor turunan kanan dan funktor turunan kiri dalam
	kerangka teori modul, yang memberikan definisi umum bagi funktor
	$\Ext$ dan $\Tor$, serta pengenalan dimensi injektif dan dimensi
	projektif suatu modul; perangkat ini segera menunjukkan kegunaannya
	dalam teori gelanggang komutatif;
	% segment-id: o014.li2.prelude.homology-history.p015
	\item penjelasan, melalui perangkat tersebut, atas berbagai konstruksi
	sebelumnya, seperti kohomologi grup serta homologi dan kohomologi
	Hochschild;
	% segment-id: o014.li2.prelude.homology-history.p016
	\item pembangunan sistematis teori umum barisan spektral dan struktur
	perkaliannya.
\end{itemize}

% segment-id: o014.li2.prelude.homology-history.p017
Istilah ``aljabar homologis'' juga berasal dari \cite{CE56}. Dengan
hadirnya teori akbar ini, aljabar homologis memasuki masa mudanya.

\subsection*{Kategori Abelian}

% segment-id: o014.li2.prelude.homology-history.p018
Dilihat dari sudut pandang geometri---lebih tepatnya, teori berkas---kerangka
teori modul belum memadai bagi aljabar homologis. Di samping kategori modul
kiri atas gelanggang $R$, yakni $R\dcate{Mod}$, kita ingin mengizinkan suatu
kategori yang lebih umum $\mathcal{A}$ tetapi tetap memiliki makna ``linear''.
Upaya pertama ke arah ini muncul dalam disertasi
doktor D.\ Buchsbaum pada tahun 1955, yang juga menjadi sebuah lampiran
dalam \cite{CE56}; ia mengabstraksikan konsep barisan eksak dari teori
modul.

% segment-id: o014.li2.prelude.homology-history.p019
Pendekatan yang kini lazim berasal dari makalah A.\ Grothendieck
\cite{Gr57}; kategori yang dihasilkan disebut kategori abelian. Kategori
ini adalah kategori aditif yang memiliki kernel dan kokernel, dengan
syarat bahwa setiap morfisme $f: M \to N$ di dalamnya bersifat ketat.
Secara kasar, keketatan tersebut bersesuaian dengan teorema isomorfisme
yang dikenal dalam teori modul,
$M/\Ker(f) \rightiso \Image(f)$. Dalam karya itu, Grothendieck juga:
\begin{itemize}
	% segment-id: o014.li2.prelude.homology-history.p020
	\item mendefinisikan funktor-$\delta$ dan funktor-$\delta$ universal
	dengan memakai konsep barisan eksak panjang, memberikan karakterisasi
	bagi yang terakhir, dan menunjukkan bahwa funktor turunan merupakan
	kasus khusus funktor-$\delta$ universal;
	% segment-id: o014.li2.prelude.homology-history.p021
	\item memberikan barisan spektral Grothendieck untuk funktor turunan
	dari komposisi funktor;
	% segment-id: o014.li2.prelude.homology-history.p022
	\item mendefinisikan suatu kelas khusus kategori abelian, yang kini
	disebut kategori Grothendieck, serta membuktikan keberadaan resolusi
	injektif di dalamnya.
\end{itemize}

% segment-id: o014.li2.prelude.homology-history.p023
Teori-teori ini memungkinkan kohomologi berkas dipelajari pada ruang yang
lebih umum dan sangat penting bagi perkembangan geometri.

\subsection*{Kategori Turunan}

% segment-id: o014.li2.prelude.homology-history.p024
Dalam disertasi doktor tahun 1963, J.-L.\ Verdier mendefinisikan kategori
turunan $\cate{D}(\mathcal{A})$ bagi suatu kategori abelian $\mathcal{A}$.
Motivasinya berasal dari kesulitan yang dihadapi
Grothendieck ketika mengkaji dualitas dalam geometri aljabar. Hambatan
tersebut mendorong pencarian kerangka baru untuk bekerja bukan lagi hanya
dengan kohomologi, melainkan dengan kompleks itu sendiri, sambil secara
formal menambahkan invers bagi kuasi-isomorfisme ke dalam kategori
kompleks $\cate{C}(\mathcal{A})$. Kuasi-isomorfisme adalah morfisme
kompleks yang menginduksi isomorfisme pada taraf kohomologi,
$\Hm^n(f): \Hm^n(X) \rightiso \Hm^n(Y)$, dengan $f: X \to Y$.
Konstruksi penambahan invers
tersebut disebut lokalisasi Gabriel--Zisman dan merupakan generalisasi
lokalisasi gelanggang pada taraf kategori.

% segment-id: o014.li2.prelude.homology-history.p025
Kategori turunan $\cate{D}(\mathcal{A})$ memuat lebih banyak informasi
daripada kohomologi, dan cara bekerja dengannya pun berbeda. Di dalam
$\cate{D}(\mathcal{A})$, kita harus memakai suatu kelas barisan morfisme
yang disebut segitiga terbedakan (\emph{distinguished triangle})
\[ X \xrightarrow{f} Y \xrightarrow{g} Z \xrightarrow{h} X[1] \]
untuk menggantikan barisan eksak pendek dalam $\mathcal{A}$ atau kategori
kompleks $\cate{C}(\mathcal{A})$. Di sini,
$X[1]$ berarti pergeseran kompleks $X$:
\[ X[1]^n = X^{n+1}, \quad d_{X[1]}^n = -d_X^n. \]
Konstruksi-konstruksi ini selanjutnya diabstraksikan oleh Verdier menjadi
suatu kelas struktur yang disebut kategori bertriangulasi
(\emph{triangulated category}). Dilihat dari
topologi, struktur bertriangulasi kategori turunan juga bersifat alami:
segitiga terbedakan dapat dianalogikan dengan barisan kofiber
(\emph{cofiber sequence}) dalam teori
homotopi. Hal ini berkaitan dengan aljabar homotopik yang akan dibahas
kelak. Atas pertimbangan teori homotopi pula, D.\ Puppe secara independen
menemukan struktur serupa pada tahun 1962, tetapi aksioma yang ia tuliskan
tidak memuat aksioma oktahedral Verdier.

% segment-id: o014.li2.prelude.homology-history.p026
Meskipun kategori turunan tidak sempurna untuk semua penerapan, bahasa ini
tetap menjadi pilihan utama para ahli geometri dan aljabar hingga awal
abad ke-21. Sebagai contoh, salah satu penerapannya yang menonjol ialah
teori modul-$\mathscr{D}$, yang terutama dikembangkan oleh Masaki Kashiwara,
Z.\ Mebkhout, dan matematikawan lain; beberapa operasi dasarnya hanya
dapat didefinisikan secara wajar pada taraf kategori turunan.

% segment-id: o014.li2.prelude.homology-history.p027
Selain itu, kategori turunan suatu gelanggang $R$,
$\cate{D}(R) := \cate{D}(R\dcate{Mod})$, dapat pula dipandang sebagai
suatu invarian bagi $R$. Jika kategori turunan dua gelanggang ekuivalen
sebagai kategori bertriangulasi, kedua gelanggang itu disebut ekuivalen
secara turunan. Berbagai persoalan dan metode seputar ekuivalensi turunan
merupakan salah satu arus utama dalam teori representasi aljabar; para
perintisnya antara lain D.\ Happel dan J.\ Rickard.

\subsection*{Aljabar Homotopik}

% segment-id: o014.li2.prelude.homology-history.p028
Pembaca tentu masih ingat bahwa asal-usul topologis aljabar homologis
adalah penguraian ruang, lalu pendefinisian kompleks rantai atau kompleks
yang sesuai untuk menyarikan invarian topologis. Penguraian pada langkah
pertama memiliki berbagai ragam; himpunan simplisial yang didefinisikan
oleh Eilenberg dan Zilber pada tahun 1950 merupakan salah satunya, dan
himpunan simplisial dapat diperumum lebih lanjut menjadi objek simplisial
dalam suatu kategori tertentu. Bahasa ini memiliki daya penjelas yang
sangat luas. Sebagai contoh, teori homotopi ruang topologis dapat
dikembangkan dalam kerangka himpunan simplisial, sedangkan konstruksi yang
disebut ``saraf'' (\emph{nerve}) menafsirkan kategori sebagai suatu jenis
khusus himpunan simplisial.

% segment-id: o014.li2.prelude.homology-history.p029
Apa yang terjadi jika sedikit ``linearitas'' ditambahkan ke dalam gambaran
ini? Korespondensi Dold--Kan (1958) memberikan hubungan langsung antara
objek simplisial dalam kategori abelian dan teori kompleks. Sebagai
contoh, untuk kategori grup abelian $\cate{Ab}$, korespondensi Dold--Kan
merupakan suatu ekuivalensi adjoin yang diwujudkan secara eksplisit oleh
sepasang funktor:
\[\begin{tikzcd}
	\mathrm{N}: \cate{sAb} \arrow[shift left, r] & \cate{Ch}_{\geq 0}(\cate{Ab}): \Gamma \arrow[shift left, l]
\end{tikzcd}\]
Di sini, $\cate{sAb}$ adalah kategori objek simplisial dalam $\cate{Ab}$,
sedangkan $\cate{Ch}_{\geq 0}(\cate{Ab})$ adalah kategori kompleks rantai
berderajat tak negatif dalam $\cate{Ab}$. Untuk kategori kompleks, kita
dapat menangani konvensi lain dengan membalik pengindeksan atau beralih ke
kategori lawan.

% segment-id: o014.li2.prelude.homology-history.p030
Korespondensi Dold--Kan memberikan tafsiran topologis bagi banyak definisi
dan konstruksi tentang kompleks rantai. Melalui metode simplisial,
suatu kelas resolusi yang secara kolektif disebut ``konstruksi bar'' dalam
aljabar homologis juga memperoleh penjelasan terpadu.

% segment-id: o014.li2.prelude.homology-history.p031
Homotopi adalah gagasan inti dalam teori simplisial. Metode aljabar yang
melibatkan objek simplisial kadang-kadang juga disebut aljabar homotopik.
Daya metode ini tampak menonjol dalam teori K tingkat tinggi yang
dikembangkan D.\ Quillen, dan penerapannya tidak terbatas pada hal itu.

% segment-id: o014.li2.prelude.homology-history.p032
Ringkasnya, metode simplisial dalam aljabar linear dapat dipandang sebagai
pertemuan, pada suatu taraf yang lebih tinggi, antara topologi---khususnya
teori homotopi---dan aljabar homologis. Teori kategori-$\infty$ merupakan
gelombang lain yang diilhami oleh pertemuan tersebut pada awal abad ke-21.
